\documentclass[10pt,letterpaper]{article}
\usepackage{cornell-notes}

\title{Clause 23.04.04.02: Non Member Comparison Operator Functions}
\author{}
\date{}
\setDocTitle{Clause 23.04.04.02: Non Member Comparison Operator Functions}
\setDocSubtitle{C++ 2024}
\setDocOwner{C++ 2024 Cornell Notes}
\setDocDate{\today}
\setCornellCollection{C++ 2024 Cornell Notes}
\setCornellUnitType{Clause}
\setCornellUnitNumber{23.04.04.02}
\setCornellUnitTitle{Non Member Comparison Operator Functions}
\hypersetup{pdftitle={Clause 23.04.04.02: Non Member Comparison Operator Functions},pdfsubject={Cornell notes},pdfauthor={}}

\begin{document}
\maketitle
\begin{CornellOverviewBox}{Overview}
\textbf{Classification:} Standard library - Normative\\[2pt]
\textbf{Learning objective:} Understand the rules, terminology, and semantic consequences associated with Non-member comparison operator functions.
\end{CornellOverviewBox}\begin{CornellSummaryBox}{Summary}
{\sffamily\bfseries\color{Accent} Summary}\par\vspace{3pt}Section 23.4.4.2 focuses on Non-member comparison operator functions. Apply the specified interface together with its mandates, effects, result conditions, exception behavior, and complexity constraints. When applying it, preserve the distinction between mandatory requirements, recommendations, permissions, and behavior deliberately left undefined, unspecified, or implementation-defined.
\end{CornellSummaryBox}

\end{document}
